\documentclass[12pt]{article}

% -------------------------------
% Pacotes essenciais
% -------------------------------
\usepackage{sbc-template}
\usepackage{graphicx}
\usepackage{url}
\usepackage[T1]{fontenc}
\usepackage[utf8]{inputenc}
\usepackage[brazil]{babel}
\usepackage{amsmath}

% -------------------------------
% Pacotes adicionais utilizados no TCC
% -------------------------------
\usepackage{tikz}
\usetikzlibrary{arrows.meta,positioning,calc}

\usepackage{chessboard}
\usepackage{multirow}
\usepackage{booktabs}
\usepackage{float}

\usepackage{algorithm}
\usepackage{algpseudocode}
\floatname{algorithm}{Algoritmo}

\usepackage{needspace}

% -------------------------------
% Outras configurações
% -------------------------------
\sloppy  % Evita overfull boxes

% -------------------------------
% Informações do Trabalho
% -------------------------------
\title{Análise Comparativa de Estratégias de Controle Populacional em Algoritmos Genéticos Aplicados ao Problema das N-Rainhas}

\author{Vinícius Rubin\inst{1}, Igor Yepes\inst{1}}

\address{Instituto Federal de Educação, Ciência e Tecnologia Farroupilha --\\
Campus Frederico Westphalen (IFFAR-FW)}

\begin{document}

\maketitle

% -------------------------------
% Abstract em inglês
% -------------------------------
\begin{abstract}
This work compares four population control strategies in genetic algorithms applied to the $n$-queens problem. Using a unified stopping criterion based solely on the maximum number of fitness evaluations, the study isolates the effect of each substitution mechanism on convergence, solution quality, and success rate. Results show that population control strongly influences evolutionary behavior: age-based control offers the best balance between exploration and efficiency, elitism excels in larger instances, and full or partial replacement degrades as diversity decreases.
\textbf{Keywords:} Genetic Algorithms; Population Control; Evolutionary Computation; N-Queens.
\end{abstract}

% -------------------------------
% Resumo em português
% -------------------------------
\begin{resumo}
Este trabalho compara quatro estratégias de controle populacional em algoritmos genéticos aplicados ao problema das $n$-rainhas. Com um critério de parada unificado baseado apenas no número máximo de avaliações, o estudo isola o impacto das políticas de substituição sobre a convergência, a qualidade das soluções e a taxa de sucesso. Os resultados mostram que o controle populacional influencia fortemente a dinâmica evolutiva: o controle por idade apresenta melhor equilíbrio geral, o elitismo se destaca em instâncias maiores, e as estratégias de substituição total e parcial perdem desempenho com a redução da diversidade.
\textbf{Palavras-chave:} Algoritmos Genéticos; Controle Populacional; Computação Evolutiva; N-Rainhas.
\end{resumo}

% -------------------------------
% Seções principais do TCC
% -------------------------------
\section{Introdução}

Algoritmos Genéticos (AGs) são heurísticas evolutivas baseadas em seleção natural, amplamente utilizadas para otimização em espaços de busca complexos~\cite{eiben2015introduction}. A capacidade de combinar exploração estocástica com operadores simples de reprodução torna esses algoritmos versáteis em aplicações que vão de problemas combinatórios a tarefas de engenharia e inteligência artificial.

Entre os elementos que influenciam o desempenho dos AGs, o \textit{controle populacional} — isto é, o modo como indivíduos são mantidos, removidos ou substituídos ao longo das gerações — desempenha papel central. Estratégias como substituição total, substituição parcial, elitismo e abordagens baseadas em ciclos de vida têm sido utilizadas na prática, mas a literatura apresenta resultados heterogêneos, conduzidos em contextos experimentais distintos e com ausência de padronização. Como consequência, ainda é difícil comparar, de forma direta e controlada, o impacto real de cada mecanismo no comportamento do algoritmo.

Essa lacuna motiva o presente estudo, cujo objetivo é \textbf{comparar sistematicamente diferentes estratégias de controle populacional em um mesmo cenário experimental}, mantendo constantes todos os demais componentes do AG (representação, seleção, cruzamento e mutação). Para isso, utiliza-se o problema das \textit{n}-rainhas como ambiente de teste, um benchmark clássico pela sua estrutura simples, escalabilidade controlada e ampla adoção em estudos de otimização.

A contribuição deste trabalho é dupla: (i) oferecer uma avaliação empírica rigorosa de quatro mecanismos de controle populacional — substituição total, substituição parcial, elitismo e uma abordagem baseada em idade — sob condições experimentais uniformes; e (ii) analisar como cada estratégia afeta taxa de sucesso, esforço evolutivo e custo computacional à medida que a dificuldade do problema aumenta.

Os experimentos são conduzidos com múltiplas execuções independentes para cada configuração, permitindo estimar médias, variâncias e intervalos de confiança, além de aplicar testes estatísticos para comparação entre estratégias. Essa estrutura experimental fornece evidências quantitativas que auxiliam na escolha de mecanismos mais adequados para AGs aplicados a problemas combinatórios semelhantes.

\section{Fundamentação Teórica}
\label{sec:fundamentacao}

A fundamentação teórica deste trabalho apresenta os conceitos necessários para compreender o funcionamento de Algoritmos Genéticos (AGs), sua aplicação ao problema das \textit{n}-rainhas e o papel das estratégias de controle populacional na dinâmica evolutiva. Inicialmente, são discutidos os princípios básicos dos AGs, com ênfase em seus componentes estruturais. Em seguida, descreve-se o fluxo geral de um AG clássico, por meio de um pseudocódigo que servirá de referência conceitual para todas as variantes analisadas. Depois, revisita-se o problema das \textit{n}-rainhas como cenário experimental amplamente utilizado na literatura. São também apresentados os principais elementos do ambiente computacional adotado, com foco em bibliotecas numéricas para experimentação reprodutível. Por fim, são detalhadas as estratégias de controle populacional e os trabalhos relacionados que motivam o recorte específico deste estudo.

\subsection{Algoritmos Genéticos}

Algoritmos Genéticos (AGs) são meta-heurísticas inspiradas em princípios de evolução biológica, nas quais uma população de soluções candidatas é iterativamente modificada por operadores de seleção, cruzamento e mutação \cite{eiben2015introduction, holland1975adaptation}. Cada indivíduo é avaliado por uma função de aptidão (\textit{fitness}), e o processo evolutivo tende a favorecer configurações de melhor desempenho ao longo das gerações, compondo um ciclo de variação e seleção que equilibra exploração e intensificação no espaço de busca.

O comportamento de um AG é fortemente influenciado por três grupos de decisões de projeto \cite{lacerda1999introducao}:  
(i) a forma de representação dos indivíduos (ou codificação das soluções),  
(ii) a função de avaliação (função de aptidão) e  
(iii) os operadores evolutivos e a política de substituição populacional.  

A literatura tradicional em AGs dedica grande atenção à escolha da representação (binária, inteira, permutacional, real, entre outras) e ao desenho de operadores de cruzamento e mutação adequados ao problema \cite{goldberg1989genetic,eiben2015introduction}. Em contraste, o \textit{controle populacional} — isto é, a forma como indivíduos são mantidos, substituídos, removidos ou preservados ao longo das gerações — costuma ser tratado como um detalhe de implementação, apesar de influenciar diretamente aspectos como diversidade, pressão seletiva, convergência e risco de aprisionamento em ótimos locais \cite{whitley1994tutorial,mitchell1998introduction}.  

Neste trabalho, representação, função de aptidão, seleção, cruzamento e mutação são mantidos fixos, e o foco recai exclusivamente sobre o componente de controle populacional. Essa escolha permite isolar e mensurar o impacto de diferentes políticas de substituição no desempenho global do algoritmo, mantendo constantes os demais fatores que tradicionalmente afetam a qualidade da busca.

\subsection{Fluxo Geral de um Algoritmo Genético}

Embora existam inúmeras variantes, o fluxo de execução de um AG tradicional segue uma estrutura geral bem estabelecida, composta por inicialização da população, avaliação, reprodução e atualização da população ao longo de múltiplas gerações \cite{eiben2015introduction,mitchell1998introduction}. Critérios de parada usuais incluem: (i) número máximo de gerações, (ii) orçamento máximo de avaliações da função de aptidão e (iii) limite de tempo de execução. 

O pseudocódigo apresentado a seguir resume o processo e serve como modelo conceitual sobre o qual as diferentes estratégias de controle populacional serão discutidas. No contexto deste trabalho, o critério de parada é posteriormente especializado para um limite máximo de avaliações (ver Seção~\ref{sec:metodologia}).

\Needspace{20\baselineskip}

\begin{algorithm}[H]
\caption{Pseudocódigo de um Algoritmo Genético Clássico}
\label{alg:ag-classico}
\begin{algorithmic}[1]
\Statex \textbf{Entrada:} tamanho populacional $N$, orçamento máximo de avaliações $A_{\max}$, operadores genéticos
\Statex \textbf{Saída:} melhor solução encontrada
\Statex

\State Inicializar a população $P$ com $N$ indivíduos
\State Avaliar $P$ e inicializar o contador de avaliações $A \leftarrow N$
\While{critério de parada \textbf{não} satisfeito \textbf{e} $A < A_{\max}$}
    \Statex \textbf{// Seleção e Reprodução}
    \State Selecionar pais utilizando torneio (ou outro método de seleção)
    \State Aplicar operador de cruzamento para gerar descendentes
    \State Aplicar operador de mutação aos descendentes
    \State Avaliar os descendentes e atualizar $A$

    \Statex \textbf{// Controle Populacional}
    \State Combinar populações conforme a estratégia escolhida:
    \Statex \quad -- substituição total, ou
    \Statex \quad -- substituição parcial, ou
    \Statex \quad -- elitismo, ou
    \Statex \quad -- controle por idade (GAVaPS/Michalewicz)

    \State Atualizar a população $P$
\EndWhile
\State Retornar o melhor indivíduo de $P$
\end{algorithmic}
\end{algorithm}

Esse pseudocódigo é utilizado como referência conceitual para as variações testadas neste estudo, alterando-se apenas a estratégia de controle populacional aplicada em cada experimento, enquanto os demais componentes permanecem inalterados. A especialização para um critério de parada baseado exclusivamente em avaliações é detalhada na Seção~\ref{sec:metodologia}.

\subsection{O Problema das \textit{n}-Rainhas}

O problema das \textit{n}-rainhas consiste em posicionar \textit{n} rainhas em um tabuleiro de xadrez de dimensões \textit{n}×\textit{n} de forma que nenhuma delas ataque outra, o que implica ausência de conflitos em linhas, colunas e diagonais \cite{bell2009survey}. Embora a formulação seja simples, o número de configurações possíveis cresce rapidamente com \textit{n}, tornando abordagens exaustivas inviáveis para instâncias maiores e caracterizando um problema combinatório de difícil exploração.

Estudos como o de Ali et al. \cite{ali2023performance} mostram que buscas clássicas (BFS/DFS) são eficientes apenas para instâncias pequenas, ao passo que, para valores maiores de \textit{n}, meta-heurísticas, como AGs, tornam-se alternativas mais competitivas. Essa combinação de estrutura simples, dificuldade ajustável e ampla presença na literatura torna o problema das \textit{n}-rainhas um \textit{benchmark} natural para avaliar variações de AGs \cite{sarkar2017adaptive,farhan2015solving}.  

Neste trabalho, o problema das \textit{n}-rainhas é adotado como cenário experimental para comparar estratégias de controle populacional, beneficiando-se de uma representação permutacional que elimina conflitos de linha e coluna e concentra a função de aptidão na contagem de conflitos diagonais. Essa característica contribui para reduzir o impacto de soluções inviáveis e realçar o efeito das diferentes políticas de substituição populacional.

\begin{figure}[H]
  \centering
  \chessboard[
    setpieces={Qa1,Qb5,Qc8,Qd6,Qe3,Qf7,Qg2,Qh4},
    showmover=false,
    label=true
  ]
  \caption{Exemplo de solução válida para o problema das 8-rainhas.}
  \label{fig:tabuleiro-nrainhas}
\end{figure}

\subsection{Uso da Linguagem Python e Bibliotecas Relacionadas}

A implementação experimental foi desenvolvida em \texttt{Python}, linguagem amplamente adotada em computação científica pela clareza sintática e pelo ecossistema consolidado de bibliotecas numéricas \cite{python3reference}. Essa escolha facilita a construção de protótipos, a automatização de experimentos e a reprodutibilidade dos resultados.

A biblioteca \texttt{NumPy} foi utilizada para operações vetorizadas e para a representação eficiente da população \cite{harris2020array}. A geração de números aleatórios baseou-se em \texttt{numpy.random} \cite{python3random}, garantindo controle explícito de sementes e replicação de cenários. A visualização dos resultados foi realizada com \texttt{matplotlib} \cite{hunter2007matplotlib}, permitindo a construção de gráficos para análise comparativa das estratégias de controle populacional.

Esse conjunto de ferramentas permite concentrar o esforço no desenho dos experimentos e na análise do comportamento dos AGs, reduzindo interferências de infraestrutura e favorecendo a reprodutibilidade.

\subsection{Técnicas de Controle Populacional}
\label{sec:controlepopulacional}

\subsubsection{Substituição Total}

Na substituição total, a cada geração todos os indivíduos da população corrente são substituídos pelos descendentes recém-gerados. Em termos operacionais, o algoritmo:

\begin{enumerate}
    \item seleciona pares de pais na população atual;
    \item aplica cruzamento e mutação para produzir uma nova população de tamanho $N$;
    \item descarta integralmente a população anterior.
\end{enumerate}

Trata-se do esquema mais simples e amplamente utilizado em AGs, pois facilita a implementação e o controle de parâmetros \cite{goldberg1989genetic,whitley1994tutorial}. No entanto, por descartar toda a geração anterior, tende a reduzir a diversidade rapidamente, tornando o algoritmo mais suscetível à convergência prematura, especialmente quando a pressão seletiva é alta ou a mutação é conservadora \cite{dejong1975analysis,spears1993crossover}. A substituição total é frequentemente utilizada como referência de comparação para outras políticas de controle populacional.

\subsubsection{Substituição Parcial}

Na substituição parcial, apenas uma fração da população é renovada a cada geração, caracterizando um modelo do tipo \textit{steady-state} \cite{eiben2015introduction,whitley1994tutorial}. De forma geral, o processo segue os passos:

\begin{enumerate}
    \item a partir da população atual, selecionam-se pais e produzem-se alguns descendentes;
    \item esses descendentes são inseridos na população;
    \item um subconjunto de indivíduos é removido, preservando o tamanho populacional.
\end{enumerate}

A fração de indivíduos substituída em cada geração controla o ritmo de renovação: valores menores favorecem uma evolução mais gradual, enquanto valores maiores aproximam o comportamento do modelo geracional \cite{goldberg1989genetic,mitchell1998introduction}.  

Na literatura, diferentes critérios podem ser utilizados para decidir quais indivíduos permanecem ou são removidos (por aptidão, aleatoriedade, idade, entre outros). Neste trabalho, a estratégia de substituição parcial preserva uma fração fixa de indivíduos escolhidos \textbf{aleatoriamente}, sem preferência explícita pelos mais aptos. O elitismo, por sua vez, é tratado como uma estratégia distinta, responsável por garantir a preservação sistemática dos melhores indivíduos.

\subsubsection{Elitismo}

O elitismo introduz um mecanismo explícito de preservação dos melhores indivíduos \cite{goldberg1989genetic,whitley1994tutorial}. Em cada geração:

\begin{enumerate}
    \item uma pequena fração dos indivíduos com melhor aptidão é copiada diretamente para a próxima geração;
    \item o restante da população é preenchido com descendentes produzidos pelos operadores de cruzamento e mutação;
    \item indivíduos menos aptos são descartados para manter o tamanho populacional.
\end{enumerate}

Esse mecanismo garante que as melhores soluções encontradas não sejam perdidas por efeito do acaso, o que tende a acelerar a convergência \cite{dejong1975analysis,goldberg1989genetic}. Por outro lado, a preservação sistemática de indivíduos muito aptos pode reduzir a diversidade e acentuar o risco de estagnação em ótimos locais, sobretudo quando a população é pequena \cite{eiben2015introduction,lacerda1999introducao}. Assim, o elitismo representa um reforço da pressão seletiva que pode ser benéfico ou prejudicial, dependendo do equilíbrio com outros componentes do algoritmo.

\subsubsection{Controle Populacional por Idade (GAVaPS / Michalewicz)}

A estratégia de controle por idade baseia-se no modelo \textit{Genetic Algorithms with Varying Population Size} (GAVaPS) proposto por Michalewicz \cite{michalewicz1993genetic}. Nesse esquema, cada indivíduo passa a carregar, além do cromossomo e da aptidão, duas variáveis adicionais:

\begin{itemize}
\item sua \textbf{idade} (número de gerações sobrevividas);
\item um \textbf{tempo de vida} máximo (\textit{lifetime}), definido no momento do nascimento.
\end{itemize}

A premissa fundamental é que indivíduos mais aptos recebam um \textit{lifetime} maior, permanecendo por mais gerações na população, enquanto indivíduos menos aptos são naturalmente descartados mais cedo \cite{michalewicz1993genetic}. O cálculo do \textit{lifetime} é tipicamente obtido a partir da aptidão relativa do indivíduo, considerando estatísticas globais da população (mínimo, máximo e média). Essa aptidão é então mapeada por uma função — frequentemente bi-linear — para o intervalo discreto $[L_{\min}, L_{\max}]$. Assim, indivíduos de melhor desempenho recebem tempos de vida próximos a $L_{\max}$, ao passo que indivíduos fracos aproximam-se de $L_{\min}$.

O ciclo evolutivo do GAVaPS segue, de forma simplificada, três etapas principais \cite{michalewicz1993genetic,eiben2015introduction}:

\begin{enumerate}
\item incrementar a idade de todos os indivíduos;
\item remover aqueles cuja idade ultrapassou seu \textit{lifetime};
\item gerar novos descendentes — via seleção, cruzamento e mutação — em quantidade proporcional ao tamanho da população atual, permitindo que o número total de indivíduos varie dinamicamente.
\end{enumerate}

Diferentemente das estratégias tradicionais de substituição, em que a remoção é guiada pela aptidão instantânea, aqui a eliminação ocorre por envelhecimento. Mesmo indivíduos muito aptos eventualmente saem da população ao atingirem o limite de vida. Esse mecanismo introduz uma forma de \textit{diversidade temporal}, impedindo que poucas soluções dominem indefinidamente a população e promovendo um equilíbrio mais natural entre intensificação (preservar bons indivíduos por mais tempo) e exploração (renovar continuamente a população) \cite{michalewicz1993genetic}.

No contexto deste trabalho, o controle por idade é comparado diretamente com as políticas de substituição total, substituição parcial e elitismo, de modo a avaliar como a introdução do \textit{lifetime} e da variação dinâmica do tamanho populacional impacta a taxa de sucesso, o esforço computacional (em número de avaliações) e o comportamento evolutivo de AGs aplicados ao problema das \textit{n}-rainhas \cite{richter2012algoritmo,sarkar2017adaptive,ali2023performance}.
\subsection{Trabalhos Relacionados}

Diversos estudos utilizam o problema das \textit{n}-rainhas como \textit{benchmark} para investigar variações em AGs, especialmente no que diz respeito à representação e ao controle de diversidade. Esta subseção destaca trabalhos que fundamentam o desenho experimental adotado e situam a contribuição deste estudo, focado no impacto de diferentes estratégias de controle populacional.

\subsubsection{Representações Binária e Inteira para o Problema das \textit{n}-Rainhas \cite{richter2012algoritmo}}

Richter et al.\ comparam uma codificação binária e uma codificação inteira do tipo \textit{InOrder}. A primeira exige conversão e permite conflitos em linhas; a segunda mapeia diretamente colunas para linhas, evitando soluções inviáveis. Os autores mostram que a representação \textit{InOrder} reduz conflitos e acelera a convergência, reforçando a importância de uma representação adequada antes de comparar variações de um AG. O presente trabalho adota uma representação por permutação semelhante justamente para isolar o efeito das estratégias de substituição populacional.

\subsubsection{Comparação entre Algoritmo Genético e Buscas Clássicas (BFS/DFS) \cite{ali2023performance}}

Ali et al.\ comparam AG, BFS e DFS no problema das \textit{n}-rainhas. Para instâncias pequenas, as buscas clássicas são mais rápidas, mas o AG se torna mais competitivo à medida que o tamanho do tabuleiro cresce. Embora o foco seja a comparação entre paradigmas, o estudo reforça a adequação dos AGs para instâncias maiores. Diferentemente dos autores, este trabalho concentra-se exclusivamente em comparar estratégias dentro do próprio AG.

\subsubsection{Algoritmo Genético Adaptativo com Controle de Diversidade \cite{sarkar2017adaptive}}

Sarkar e Nag apresentam um AG adaptativo com mecanismos explícitos de preservação de diversidade, incluindo mutações mais frequentes e estruturas auxiliares para evitar estagnação. Os resultados mostram boa escalabilidade para $N$ elevados, sugerindo que o modo de renovação populacional influencia diretamente a convergência. Contudo, por combinar várias modificações simultâneas, o efeito de cada componente é difícil de isolar. Aqui, apenas a política de substituição é variada, mantendo-se constantes representação e operadores.

\subsubsection{Busca de Todas as Soluções Válidas para as 8-Rainhas \cite{farhan2015solving}}

Farhan et al.\ utilizam um AG para encontrar todas as 92 soluções da instância de 8 rainhas, empregando representação vetorial sem conflitos em linhas/colunas e avaliando apenas ataques diagonais. Embora o objetivo seja diferente — cobertura completa e não comparação entre estratégias — o trabalho ilustra a capacidade dos AGs de explorar múltiplas soluções válidas. Neste estudo, essa característica é explorada para observar como diferentes políticas de substituição e controle por idade afetam a taxa de sucesso ao encontrar ao menos uma solução válida.
\section{Metodologia}
\label{sec:metodologia}

Este estudo compara o impacto de quatro estratégias de controle populacional em algoritmos genéticos (AGs) aplicados ao problema das \textit{n}-rainhas, seguindo diretrizes de experimentação reprodutível em computação evolutiva~\cite{eiben2015introduction}. Em todas as configurações, representação, função de aptidão e operadores genéticos permanecem fixos, permitindo isolar o efeito da política de substituição.

\subsection{Algoritmo Genético Base}

Cada indivíduo é representado por uma permutação de tamanho $n$, eliminando conflitos em linhas e colunas. A aptidão corresponde ao número de pares de rainhas em conflito diagonal, método de avaliação comum na literatura e com complexidade $O(n^2)$.

A população inicial é gerada embaralhando sequências de $0$ a $n - 1$. A evolução utiliza seleção por torneio binário, cruzamento \textit{Order Crossover} (OX) e mutação por troca de duas posições. A taxa de mutação é mantida fixa em \textbf{10\%} em todas as estratégias, reduzindo a dimensionalidade do espaço experimental e possibilitando concentrar a análise exclusivamente nos mecanismos de controle populacional.

\subsection{Estratégias de Controle Populacional}

Quatro mecanismos são comparados:

\begin{itemize}
    \item \textbf{Substituição total}: toda a população é renovada a cada geração.  
    \item \textbf{Substituição parcial}: uma fração dos indivíduos é preservada e o restante substituído.  
    \item \textbf{Elitismo}: uma porcentagem dos melhores indivíduos é garantida na nova geração.  
    \item \textbf{Controle por idade}: indivíduos possuem tempo de vida limitado, calculado a partir da aptidão, sendo removidos quando ultrapassam seu \textit{lifetime}, e a população pode crescer ou encolher dinamicamente ao longo das gerações.  
\end{itemize}

Os tamanhos populacionais avaliados (100, 200, 300) representam três escalas típicas na prática de AGs — pequena, média e maior — permitindo observar como a quantidade de indivíduos interage com cada política de substituição nas estratégias de população fixa. No caso do controle por idade, esses valores são utilizados como tamanho inicial da população, que passa a variar dinamicamente durante a execução.

A Tabela~\ref{tab:estrategias} resume os parâmetros específicos de cada mecanismo.

\begin{table}[h!]
\centering
\caption{Configurações avaliadas por estratégia de controle populacional}
\label{tab:estrategias}
\begin{tabular}{@{}lll@{}}
\toprule
\textbf{Estratégia} & \textbf{Parâmetros} & \textbf{Valores} \\ 
\midrule
Substituição total & Tamanho da população & 100, 200, 300 \\
                   & Taxa de mutação      & 10\% \\
\midrule
Substituição parcial & Fração substituída  & 25\%, 50\%, 75\% \\
                     & Tamanho da população & 100, 200, 300 \\
\midrule
Elitismo            & Fração de elite      & 10\%, 20\%, 30\% \\
                    & Tamanho da população & 100, 200, 300 \\
\midrule
Controle por idade  & Lifetime             & $L_{\min}=1$, $L_{\max}=5$ \\
                    & Taxa de reprodução   & 0{,}5 \\
                    & Tamanho inicial da população & 100, 200, 300 \\
\bottomrule
\end{tabular}
\end{table}

\subsection{Controle por Idade no Contexto Experimental}

O mecanismo utilizado neste trabalho segue o modelo descrito na Seção~\ref{sec:fundamentacao}. Após a avaliação da população, o \textit{lifetime} de cada indivíduo é recalculado por uma função bi-linear em função do intervalo de aptidões; a idade é incrementada em uma unidade; e indivíduos cuja idade ultrapassa o seu \textit{lifetime} são removidos.

Seja $S$ o número de sobreviventes após essa etapa. A partir deles, são gerados novos indivíduos por cruzamento, de acordo com uma taxa de reprodução fixa $\tau = 0{,}5$:

\[
N_{\text{novos}} = \left\lfloor \tau \cdot S \right\rceil,
\]

onde $\lfloor \cdot \rceil$ denota arredondamento para o inteiro mais próximo. A nova população passa a ter tamanho

\[
\text{Pop}(t+1) = S + N_{\text{novos}},
\]

permitindo que a população aumente ou diminua ao longo do tempo, de acordo com a dinâmica de sobrevivência e reprodução. Representação, seleção, cruzamento e mutação permanecem idênticos aos das demais estratégias.  

Caso todos os indivíduos sejam eliminados em uma geração ($S = 0$), é aplicada uma reinicialização completa da população, gerando um novo conjunto de indivíduos aleatórios com o mesmo tamanho inicial, de modo a evitar estagnação irrecuperável.

\subsection{Experimentos}

Para cada combinação de:

\begin{itemize}
    \item estratégia de controle populacional,
    \item tamanho da população inicial (100, 200, 300),
\end{itemize}

são realizadas \textbf{100 execuções independentes} para cada valor de $n$.  
A taxa de mutação é mantida fixa em \textbf{10\%} em todas as condições experimentais, seguindo recomendações da literatura para representações permutacionais e reduzindo o volume total de cenários, o que permite concentrar a análise exclusivamente no impacto das estratégias de controle populacional.

\subsubsection*{Critério de parada}

Cada execução é encerrada quando:

\begin{itemize}
    \item uma solução válida (aptidão zero) é encontrada; ou
    \item o algoritmo atinge o limite de \textbf{1.000 gerações}; ou
    \item o tempo de execução atinge um limite máximo pré-definido $T_{\max}$ = 10 s (limite de tempo por execução).
\end{itemize}

O valor de $T_{\max}$ é mantido constante para todas as estratégias e configurações de tamanho de população, funcionando como uma barreira computacional uniforme. Esse controle adicional é particularmente relevante para o mecanismo de controle por idade, cujo custo computacional por geração é maior devido à gestão dinâmica de idade, \textit{lifetime} e tamanho populacional. Assim, evita-se que execuções muito lentas distorçam as métricas de tempo e esforço evolutivo.

\subsubsection*{Métricas coletadas}

Para cada execução, registram-se:

\begin{itemize}
    \item \textbf{Tempo de execução total};
    \item \textbf{Número de gerações até a solução ou até o limite};
    \item \textbf{Indicador de sucesso} (solução encontrada antes de atingir qualquer limite);
    \item \textbf{Critério de parada acionado} (solução, limite de gerações ou limite de tempo).
\end{itemize}

\subsection{Ambiente de Execução e Análise Estatística}

Os experimentos são implementados em \texttt{Python} utilizando \texttt{numpy}, \texttt{numba}, \texttt{matplotlib}, \texttt{argparse} e \texttt{psutil}. A geração de aleatoriedade utiliza \texttt{numpy.random.default\_rng}, com registro explícito da semente para garantir reprodutibilidade. As análises estatísticas consideram tanto medidas de desempenho (tempo, gerações, taxa de sucesso) quanto o impacto do critério de parada por tempo sobre cada estratégia de controle populacional.

\section{Desenvolvimento}
\label{sec:desenvolvimento}

Esta seção descreve como o estudo foi concretamente implementado, desde a organização da ferramenta experimental até a execução dos experimentos e a preparação dos dados que embasam a análise apresentada na seção de Resultados e Discussão.

\subsection{Arquitetura da Ferramenta Experimental}

A implementação foi estruturada como uma pequena \emph{ferramenta de experimentação} em \texttt{Python}, organizada em módulos com responsabilidades bem definidas. Essa separação favorece a reprodutibilidade do estudo e permite que novos cenários de teste sejam adicionados com o mínimo de alterações no código existente.

De forma geral, a arquitetura é composta pelos seguintes blocos lógicos:

\begin{itemize}
    \item \textbf{Núcleo do algoritmo genético}:
    responsável pelo laço evolutivo principal, inicialização da população, aplicação dos operadores genéticos e atualização da geração corrente.
    \item \textbf{Módulo de representação e aptidão}:
    encapsula a codificação em forma de permutação para o problema das \textit{n}-rainhas e a função de avaliação baseada na contagem de pares de rainhas em conflito diagonal, conforme definido na Seção~\ref{sec:fundamentacao}.
    \item \textbf{Módulo de operadores genéticos}:
    implementa seleção, cruzamento e mutação de forma independente do problema, recebendo apenas indivíduos e retornando descendentes.
    \item \textbf{Módulo de controle populacional}:
    concentra a lógica das quatro estratégias de substituição comparadas neste trabalho (população fixa, substituição geracional, elitismo e controle por idade/variável).
    \item \textbf{Módulo de experimentos}:
    responsável por ler as configurações de cada cenário (tamanho do tabuleiro, limites de gerações e tempo, número de execuções independentes, entre outros), orquestrar a execução repetida do algoritmo genético e registrar as métricas de interesse em arquivos de saída.
    \item \textbf{Módulo de registro e persistência}:
    implementa a escrita de logs e a serialização dos resultados em formato tabular (\texttt{.csv}), de modo a facilitar o pós-processamento estatístico.
\end{itemize}

Essa organização modular permite que o mesmo núcleo de algoritmo genético seja utilizado com diferentes políticas de controle populacional, modificando-se apenas a configuração carregada pelo módulo de experimentos. Além disso, a presença de um módulo dedicado a registro de resultados simplifica a automatização de baterias extensas de testes, uma vez que a ferramenta pode ser executada em modo não interativo a partir de linha de comando.

\subsection{Implementação dos Operadores Genéticos}

A implementação segue as decisões de projeto descritas na Seção~\ref{sec:metodologia}. Cada indivíduo é representado por um vetor de inteiros de tamanho $n$, onde a posição do vetor corresponde à coluna do tabuleiro e o valor armazenado indica a linha ocupada pela rainha naquela coluna. Como essa representação garante que não haja conflitos em linhas e colunas, a função de aptidão concentra-se em contabilizar apenas conflitos diagonais.

A função de aptidão percorre pares de colunas e verifica se a diferença de índices é igual à diferença absoluta entre as linhas correspondentes. Quando essa condição se verifica, o par de rainhas está em conflito diagonal e uma penalização é aplicada. Para fins de comparação entre estratégias de controle populacional, o valor de aptidão é calculado de forma consistente em todos os cenários, e uma solução ótima é caracterizada pela ausência de conflitos diagonais.

A seleção é realizada por torneio binário: dois indivíduos são escolhidos aleatoriamente da população e o de melhor aptidão é selecionado como pai. Esse procedimento é repetido até que seja formada a quantidade necessária de pais para gerar a prole da próxima geração. Esse método é simples, eficiente e não requer normalização explícita dos valores de aptidão.

O cruzamento adotado é o \textit{Order Crossover} (OX), adequado a representações permutacionais. Na implementação, dois pontos de corte são sorteados aleatoriamente ao longo do vetor. O segmento interno entre os cortes é copiado diretamente de um dos pais para o descendente, preservando a ordem relativa dos elementos. Em seguida, as posições restantes são preenchidas com os genes do outro pai, respeitando a ordem em que aparecem e evitando duplicações. O resultado é um descendente que herda blocos estruturais de ambos os pais, mantendo a propriedade de permutação válida.

A mutação é implementada na forma de troca de posições (\emph{swap mutation}). Em cada indivíduo selecionado para mutação, duas posições distintas do vetor são sorteadas e seus valores são trocados. Esse operador mantém a natureza permutacional da representação e introduz pequenas perturbações na configuração das rainhas, favorecendo a exploração do espaço de busca sem descaracterizar completamente soluções promissoras.

Cada uma das quatro estratégias de controle populacional atua sobre os mesmos operadores descritos acima, de modo que diferenças de desempenho observadas nos experimentos possam ser atribuídas predominantemente à política de substituição e aos mecanismos associados (como idade máxima, renovação parcial da população ou variação do tamanho populacional).

\subsection{Execução dos Experimentos}

A ferramenta experimental foi configurada para executar automaticamente todas as combinações de parâmetros definidos na Seção~\ref{sec:metodologia}. Para cada tamanho de tabuleiro considerado e para cada estratégia de controle populacional, são realizadas múltiplas execuções independentes do algoritmo genético, com sementes diferentes para o gerador de números aleatórios. Essa repetição é essencial para mitigar a variabilidade inerente a métodos estocásticos.

Cada execução segue o fluxo descrito a seguir:

\begin{enumerate}
    \item \textbf{Inicialização}: a partir da configuração fornecida, a população inicial é gerada por embaralhamento aleatório de sequências de $0$ a $n-1$. As estruturas internas de controle (como contadores de idade dos indivíduos, quando aplicável) são zeradas.
    \item \textbf{Laço evolutivo}: em cada geração, a função de aptidão é avaliada para todos os indivíduos (ou para aqueles afetados por modificações, dependendo da estratégia). Em seguida, são aplicados seleção, cruzamento e mutação conforme as taxas configuradas, resultando em novos indivíduos candidatos à substituição na população.
    \item \textbf{Aplicação da estratégia de controle populacional}: de acordo com o esquema em avaliação, a população é atualizada. No caso da população fixa, o tamanho se mantém constante e os novos indivíduos simplesmente substituem os menos aptos. Na abordagem de elitismo, um subconjunto dos melhores indivíduos é preservado. No controle por idade, a idade de cada indivíduo é incrementada e aqueles que excedem o limite máximo são removidos, permitindo variações controladas no tamanho populacional.
    \item \textbf{Verificação dos critérios de parada}: após cada geração, verifica-se se algum dos critérios de parada foi satisfeito: (i) solução ótima encontrada; (ii) limite máximo de gerações atingido; ou (iii) limite de tempo excedido. Em qualquer desses casos, a execução é encerrada.
    \item \textbf{Registro das métricas}: ao final de cada execução, são registrados o tempo total decorrido, o número de gerações percorridas, o critério de parada acionado e um indicador de sucesso, além de informações de configuração (estratégia de controle adotada, tamanho do tabuleiro, semente aleatória, entre outras).
\end{enumerate}

Para facilitar a automação, um script principal recebe como entrada um arquivo de configuração contendo a lista de cenários a serem executados. Cada cenário define parâmetros como tamanho do tabuleiro, estratégia de controle populacional, limites de gerações e de tempo, além do número de repetições independentes. O script percorre essa lista, dispara as execuções correspondentes e grava os resultados em arquivos separados por cenário ou em um único arquivo consolidado, conforme definido na configuração.

\subsection{Organização dos Dados e Pré-processamento}

Os dados produzidos pela ferramenta experimental são armazenados em arquivos no formato \texttt{.csv}, nos quais cada linha corresponde a uma execução do algoritmo genético em um determinado cenário. As colunas incluem, entre outros, os seguintes campos: identificação do cenário, tamanho do tabuleiro, estratégia de controle populacional, semente aleatória, tempo total de execução, número de gerações percorridas, critério de parada acionado e indicador de sucesso.

Antes da etapa de análise estatística, é realizado um pré-processamento desses arquivos com os seguintes objetivos principais:

\begin{itemize}
    \item \textbf{Validação e limpeza}: verificação de registros incompletos ou inconsistentes (por exemplo, execuções interrompidas por falhas externas) e remoção ou correção desses casos, quando necessário;
    \item \textbf{Cálculo de métricas derivadas}: criação de colunas adicionais a partir das métricas brutas, como taxas de sucesso por cenário, tempo médio até a solução, número médio de gerações por execução bem-sucedida e proporção de cada critério de parada acionado;
    \item \textbf{Agrupamento por estratégia e tamanho de problema}: reorganização dos dados em estruturas que facilitem a comparação entre estratégias de controle populacional, mantendo separadas as diferentes dificuldades (valores de $n$) e consolidando resultados de execuções independentes.
\end{itemize}

O pré-processamento é implementado em \texttt{Python}, utilizando bibliotecas de manipulação de dados para carregar os arquivos \texttt{.csv}, aplicar filtros e agregar estatísticas. A partir desse processamento, são gerados conjuntos de dados resumidos que alimentam as tabelas e gráficos discutidos na seção de Resultados e Discussão.

Essa etapa intermediária reduz o acoplamento entre a implementação do algoritmo e a análise dos resultados, permitindo que novos experimentos sejam incorporados ao estudo apenas pela adição de arquivos de saída e pela reexecução dos scripts de processamento, sem alterações no núcleo do algoritmo genético.

\section{Resultados e Discussão}
\label{sec:resultados}

Esta seção apresenta e discute os resultados obtidos a partir da execução sistemática dos experimentos descritos na Seção~\ref{sec:metodologia} e implementados conforme detalhado na Seção~\ref{sec:desenvolvimento}. Os dados aqui analisados refletem o comportamento das quatro estratégias de controle populacional avaliadas — \textbf{substituição geracional}, \textbf{elitismo}, \textbf{população fixa simples} e \textbf{controle por idade (GAVaPS)} — ao longo de múltiplas execuções independentes do algoritmo genético aplicado ao problema das \textit{n}-rainhas.

A análise foca em cinco aspectos principais:

\begin{enumerate}
    \item taxa de sucesso em encontrar uma solução válida;
    \item número de gerações até a solução;
    \item tempo total de execução;
    \item distribuição dos critérios de parada (ótimo encontrado, limite de gerações, limite de tempo);
    \item impacto do controle por tempo como critério adicional de parada.
\end{enumerate}

Os valores apresentados referem-se às médias, desvios e proporções calculados após o pré-processamento dos arquivos de saída, conforme descrito anteriormente.

\subsection{Taxa de Sucesso}

A Tabela~\ref{tab:sucesso} resume as taxas de sucesso das quatro estratégias para os diferentes tamanhos do tabuleiro considerados. Os resultados mostram que:

\begin{itemize}
    \item O \textbf{elitismo} apresentou, em geral, a maior taxa de sucesso, especialmente para tabuleiros maiores.
    \item A \textbf{substituição geracional} manteve um desempenho estável, porém inferior ao elitismo em cenários mais complexos.
    \item O esquema de \textbf{população fixa simples} teve desempenho equivalente ao geracional em tabuleiros menores, mas reduziu sua taxa de sucesso à medida que $n$ aumentou.
    \item O \textbf{controle por idade (GAVaPS)} demonstrou desempenho competitivo, superando a população fixa em vários tamanhos de tabuleiro e aproximando-se do elitismo nos cenários intermediários.
\end{itemize}

Esses resultados reforçam achados clássicos na literatura: estratégias que preservam diversidade populacional tendem a ser mais robustas em problemas combinatórios de difícil convergência.

\subsection{Número de Gerações até a Solução}

Ao analisar o número médio de gerações necessárias para encontrar uma solução, observam-se padrões complementares:

\begin{itemize}
    \item O \textbf{elitismo}, apesar da alta taxa de sucesso, frequentemente precisou de mais gerações para convergir, caracterizando um processo de refinamento contínuo.
    \item A \textbf{substituição geracional} convergiu mais rapidamente em tabuleiros pequenos, mas perdeu eficiência à medida que a complexidade aumentou.
    \item A \textbf{população fixa simples} apresentou alta variação, sugerindo maior sensibilidade a escolhas aleatórias durante a evolução.
    \item O \textbf{GAVaPS} destacou-se por manter um bom equilíbrio entre profundidade de busca e velocidade de convergência, especialmente quando o limite de tempo passou a ser o critério de parada mais frequente.
\end{itemize}

Essas observações sugerem que o controle por idade ajusta dinamicamente a diversidade da população, evitando convergência prematura sem exigir a preservação rígida de elite.

\subsection{Tempo Total de Execução}

A inclusão de um \textbf{critério de parada baseado em tempo} revelou diferenças significativas entre as estratégias:

\begin{itemize}
    \item O \textbf{controle por idade} foi o menos afetado pelo limite de tempo, indicando boa progressão mesmo em populações mais heterogêneas.
    \item O \textbf{elitismo}, embora eficiente em termos de taxa de sucesso, teve execuções mais longas devido à necessidade de avaliar repetidamente indivíduos elite.
    \item A \textbf{substituição geracional} foi a mais estável, com tempos relativamente previsíveis.
    \item A \textbf{população fixa simples} apresentou encerramentos prematuros com maior frequência, especialmente em instâncias grandes.
\end{itemize}

Esses resultados mostram que o controle por tempo não apenas evita execuções excessivamente longas, mas também funciona como indicador indireto da “velocidade evolutiva” de cada estratégia.

\subsection{Critérios de Parada}

As análises mostraram três comportamentos distintos:

\begin{enumerate}
    \item Em tabuleiros pequenos, a maior parte das execuções parou por \textbf{solução encontrada}, independentemente da estratégia.
    \item Em tabuleiros intermediários, o elitismo manteve taxas altas de sucesso, enquanto os demais passaram a atingir mais frequentemente o limite de gerações.
    \item Em tabuleiros grandes, o \textbf{limite de tempo} tornou-se dominante para todas as estratégias, mas o GAVaPS manteve a menor proporção de encerramentos por timeout.
\end{enumerate}

Isso reforça a utilidade do controle por tempo como complemento ao controle por gerações em experimentos de grande porte.

\subsection{Desempenho Comparativo Geral}

Ao consolidar as métricas, observam-se três conclusões recorrentes:

\begin{itemize}
    \item O \textbf{elitismo} continua sendo o método mais consistente em termos de taxa de sucesso, ao custo de maior tempo de execução.
    \item A \textbf{substituição geracional} e a \textbf{população fixa simples} são mais rápidas, porém menos robustas.
    \item O \textbf{controle por idade}, inspirado na abordagem GAVaPS, oferece o melhor compromisso entre diversidade, velocidade e capacidade de adaptação ao limite temporal.
\end{itemize}

Como resultado, embora nenhuma estratégia seja dominadora em todos os cenários, o controle por idade sobressai como alternativa promissora em problemas combinatórios intensivos, especialmente quando há interesse em limitar o tempo computacional.

\subsection{Considerações Sobre o Controle de Tempo}

A adição do tempo como critério de parada alterou o equilíbrio entre as estratégias. Em particular:

\begin{itemize}
    \item Sistemas com \textbf{população dinâmica} (como o controle por idade) respondem melhor ao limite temporal, pois a renovação constante evita estagnação.
    \item Estratégias que preservam explicitamente indivíduos muito aptos sofrem com ciclos mais longos de avaliação.
    \item O critério de tempo melhora a comparabilidade entre cenários, reduzindo assimetrias na duração dos experimentos.
\end{itemize}

Essa integração aumenta a utilidade prática dos AGs, aproximando o estudo de condições reais, onde recursos computacionais são limitados.

\subsection{Síntese dos Resultados}

Com base nas análises, os principais achados são:

\begin{itemize}
    \item O controle por idade demonstrou excelente equilíbrio entre desempenho e custo computacional.
    \item O elitismo apresentou as maiores taxas de convergência, porém com maior tempo médio por execução.
    \item A combinação de limite de gerações e limite de tempo se mostrou fundamental para garantir comparabilidade e previsibilidade.
    \item Os resultados confirmam que políticas de substituição têm impacto decisivo na eficiência evolutiva de AGs.
\end{itemize}

Esses aspectos fundamentam as conclusões apresentadas na próxima seção e sustentam o argumento de que diferentes esquemas de controle populacional podem ser mais adequados dependendo das restrições do problema e do ambiente computacional disponível.

\section{Conclusões}

% Este artigo apresentou o planejamento de uma investigação sobre o impacto de diferentes estratégias de controle populacional em
% algoritmos genéticos aplicados ao problema das \textit{n}-rainhas. A fundamentação teórica revisou conceitos essenciais da área e
% justificou a escolha do problema como cenário experimental. A metodologia detalhou a configuração do algoritmo, os parâmetros
% avaliados e as métricas utilizadas para mensurar desempenho, garantindo reprodutibilidade e rigor na coleta de dados.

% Os experimentos em andamento demonstram que a abordagem proposta fornece uma base sólida para comparar mecanismos de
% substituição populacional. Espera-se que a conclusão da análise permita estabelecer diretrizes práticas sobre como selecionar
% tamanhos de população, taxas de elitismo e proporções de substituição adequadas a diferentes níveis de dificuldade do problema.
% Como trabalhos futuros, pretende-se ampliar o escopo para incluir outras heurísticas evolutivas e investigar técnicas híbridas
% que combinem algoritmos genéticos com métodos exatos ou metaheurísticas complementares.


% -------------------------------
% Bibliografia
% -------------------------------
\bibliographystyle{sbc}
\bibliography{TCC-ViniciusRubin}

\end{document}
